\clearpage
\section{The actual enlarged rank-two inducing carrier}\label{app:rank2}

Status: ordinary proof. The construction below supplies an actual
arithmetic finite Heisenberg pair and exact recovery of an already fixed
representation. It does not change the previously audited
zero-dimensional BK quotient at rank two. No Lean verification or H5
comparison is claimed.

\hypertarget{rank2-exact-setting-and-result}{%
\subsection{1. Exact setting and
result}\label{rank2-exact-setting-and-result}}

Let F=F\_2((t)) and G=GL2(F). Let {[}A,q,0,beta{]} be a minimal simple
stratum with K=F{[}beta{]} totally ramified of degree2 and A its
complete period-two hereditary order. Assume

\begin{Verbatim}[fontsize=\small,breaklines=true,breakanywhere=true]
q=2m+1>=3,       m>=1.
\end{Verbatim}

Write R=rad(A), \(U_A^a=1+R^a\) for a\textgreater=1, and
\(U_K^a=1+\mathfrak p_K^a\). Let theta be a full simple character.
Since q is odd and the residue field is F2, the actual maximal simple
type group and character groups coincide:

\begin{Verbatim}[fontsize=\small,breaklines=true,breakanywhere=true]
J=J^1=H^1=U_K^1 U_A^(m+1).
\end{Verbatim}

Fix the full actual extension Lambda of theta to

\begin{Verbatim}[fontsize=\small,breaklines=true,breakanywhere=true]
N=K×J,
\end{Verbatim}

and the already fixed irreducible supercuspidal candidate

\begin{Verbatim}[fontsize=\small,breaklines=true,breakanywhere=true]
Pi=c-Ind_N^G Lambda.
\end{Verbatim}

Lambda is one-dimensional here. Its value on a uniformizer of K,
including the noncompact scalar normalization, is part of the fixed
data.

Define the larger compact open arithmetic subgroup

\begin{Verbatim}[fontsize=\small,breaklines=true,breakanywhere=true]
P=U_K^1 U_A^m.
\end{Verbatim}

Then the following statements hold.

\begin{enumerate}
\def\labelenumi{\arabic{enumi}.}
\item
  J is normal of index2 in P, and P∩N=J.
\item
  E=Ind\_J\^{}P theta is an actual irreducible two-dimensional
  P-representation.
\item
  There is an explicitly defined normal subgroup H\_new⊂J of index2 such
  that P/H\_new≅F2² and theta\_new=theta\textbar H\_new has
  nondegenerate commutator pairing on that quotient. E is its Heisenberg
  representation.
\item
  E embeds into the restriction of the same Pi, as the subspace of
  inducing functions supported on NP.
\item
  With X=c-Ind\_J\^{}G theta, a uniformizer varpi\_K and
  s=Lambda(varpi\_K), there is an exact quotient

\begin{Verbatim}[fontsize=\small,breaklines=true,breakanywhere=true]
Pi ≅ X/(T−s)X,
(Tf)(g)=f(varpi_K g).
\end{Verbatim}

  Induction in stages identifies X with c-Ind\_P\^{}G E. This gives an
  exact same-candidate recovery using the retained noncompact scalar,
  not merely a nonzero cyclic-generation assertion.
\end{enumerate}

The new Heisenberg pair is an actual pair of compact subgroups of
GL2(F), but is not asserted to be the BK \(H^1/J^1\) pair of
a maximal simple type. Whether the original H2 realization relation
allows this enlargement is a separate obligation.

\hypertarget{rank2-the-compact-enlargement-is-a-group-with-normal-subgroup-of-index-two}{%
\subsection{2. The compact enlargement is a group, with normal subgroup
of index
two}\label{rank2-the-compact-enlargement-is-a-group-with-normal-subgroup-of-index-two}}

K× normalizes A and every power of its radical, so U\_K\^{}1 normalizes
U\_A\^{}m. Therefore P is a subgroup and is compact open. It is
contained in U\_A\^{}1 because m\textgreater=1 and
\(U_K^1\subset U_A^1\).

The standard radical calculation gives

\begin{Verbatim}[fontsize=\small,breaklines=true,breakanywhere=true]
[U_A^a,U_A^b]⊂U_A^(a+b),       a,b>=1.
\end{Verbatim}

Indeed, for x∈R\^{}a and y∈R\^{}b, the difference (1+x)(1+y)−(1+y)(1+x)
is in R\^{}(a+b), and multiplication by unit inverses preserves that
ideal.

Consequently

\begin{Verbatim}[fontsize=\small,breaklines=true,breakanywhere=true]
[U_A^m,U_K^1]⊂U_A^(m+1),
[U_A^m,U_A^(m+1)]⊂U_A^(2m+1)⊂U_A^(m+1).
\end{Verbatim}

These facts, and normalization by U\_K\^{}1, show that J is normal in P.
They also show

\begin{Verbatim}[fontsize=\small,breaklines=true,breakanywhere=true]
[P,J]⊂U_A^(m+1).                                      (1)
\end{Verbatim}

To calculate the index, use

\begin{Verbatim}[fontsize=\small,breaklines=true,breakanywhere=true]
U_A^m∩J=U_K^m U_A^(m+1).
\end{Verbatim}

For if ku∈U\_A\^{}m with k∈U\_K\^{}1 and u∈U\_A\^{}(m+1), then
\(k\in K^\times\cap U_A^m=U_K^m\). The reverse inclusion is immediate.
The complete period-two A-grade \(R^m/R^{m+1}\) has
F2-dimension2, and its field grade \(\mathfrak p_K^m/\mathfrak p_K^{m+1}\) has
dimension1. Thus

\begin{Verbatim}[fontsize=\small,breaklines=true,breakanywhere=true]
[P:J]=[U_A^m:U_K^m U_A^(m+1)]=2.                      (2)
\end{Verbatim}

Finally, if kj∈P with k∈K× and j∈J, then \(k\in Pj^{-1}\subset U_A^1\).
Hence k∈U\_K\^{}1 and kj∈J. Therefore

\begin{Verbatim}[fontsize=\small,breaklines=true,breakanywhere=true]
P∩N=J.                                                (3)
\end{Verbatim}

These arguments use the actual group intersections. They do not replace
P or J by a vector-space surrogate.

\hypertarget{rank2-irreducibility-and-the-conjugate-character-ratio}{%
\subsection{3. Irreducibility and the conjugate-character
ratio}\label{rank2-irreducibility-and-the-conjugate-character-ratio}}

For the minimal simple character, the primary BK intertwining theorem
gives

\begin{Verbatim}[fontsize=\small,breaklines=true,breakanywhere=true]
I_G(theta)=J K×J=N.
\end{Verbatim}

Here K× normalizes J and theta. In particular, for g∈P\textbackslash J,
(3) says g does not intertwine theta. Because J is normal in P, this
means that the two characters

\begin{Verbatim}[fontsize=\small,breaklines=true,breakanywhere=true]
theta,
theta^g(j)=theta(gjg^−1)
\end{Verbatim}

are distinct. The index-two induction E=Ind\_J\^{}P theta is therefore
irreducible. This also follows directly from the two distinct J-weight
lines, which are interchanged by g.

Define the actual character of J

\begin{Verbatim}[fontsize=\small,breaklines=true,breakanywhere=true]
lambda(j)=theta^g(j)/theta(j).
\end{Verbatim}

It is nontrivial. By (1),

\begin{Verbatim}[fontsize=\small,breaklines=true,breakanywhere=true]
lambda(j)=theta(gjg^−1j^−1),
\end{Verbatim}

and the argument of theta lies in U\_A\^{}(m+1). The full simple
character has exact restriction

\begin{Verbatim}[fontsize=\small,breaklines=true,breakanywhere=true]
theta(1+x)=psi_F(tr_F(beta x)),       x∈R^(m+1).
\end{Verbatim}

Since F has characteristic2, its additive characters take values in
\{1,−1\}. Consequently

\begin{Verbatim}[fontsize=\small,breaklines=true,breakanywhere=true]
lambda(J)={1,−1}.                                    (4)
\end{Verbatim}

No assertion about shallow values of theta is used here. Theta may have
a larger finite image on J. Only its ratio with its actual conjugate is
proved to have order2.

As g²∈J, conjugation by g² fixes the character theta. Thus
\(\lambda^g=\lambda^{-1}=\lambda\). It follows that

\begin{Verbatim}[fontsize=\small,breaklines=true,breakanywhere=true]
H_new=ker(lambda)
\end{Verbatim}

is normal in P. It has index2 in J and index4 in P.

\hypertarget{rank2-the-actual-finite-heisenberg-pair}{%
\subsection{4. The actual finite Heisenberg
pair}\label{rank2-the-actual-finite-heisenberg-pair}}

Use the function model of E with f(jp)=theta(j)f(p) and right
translation by P. Coordinates are f(1),f(g). Then for j∈J,

\begin{Verbatim}[fontsize=\small,breaklines=true,breakanywhere=true]
E(j)=diag(theta(j),theta^g(j))
    =theta(j) diag(1,lambda(j)),
\end{Verbatim}

while

\begin{Verbatim}[fontsize=\small,breaklines=true,breakanywhere=true]
E(g)=[[0,1],[theta(g²),0]].                            (5)
\end{Verbatim}

Therefore E(g)²=theta(g²)I. Choose j\_0∈J with lambda(j\_0)=−1, and put

\begin{Verbatim}[fontsize=\small,breaklines=true,breakanywhere=true]
Z=theta(j_0)^−1 E(j_0)=diag(1,−1),
X=E(g).
\end{Verbatim}

The exact matrix identity is

\begin{Verbatim}[fontsize=\small,breaklines=true,breakanywhere=true]
XZX^−1=−Z.                                           (6)
\end{Verbatim}

The projective image of J has order2, generated by Z, while X is
off-diagonal and contributes an independent projective involution. Thus
the projective image of P is exactly C2×C2, with nontrivial commutator
(6). Its projective kernel is precisely H\_new. In particular

\begin{Verbatim}[fontsize=\small,breaklines=true,breakanywhere=true]
P/H_new≅F2².
\end{Verbatim}

On H\_new, E is the scalar theta\_new=theta\textbar H\_new. This
character is P-invariant: conjugation by J fixes theta and conjugation
by g changes theta by lambda, which is trivial on H\_new. The resulting
alternating form

\begin{Verbatim}[fontsize=\small,breaklines=true,breakanywhere=true]
(pH_new,p'H_new)↦theta_new([p,p'])
\end{Verbatim}

is nondegenerate, since its two displayed generators have commutator −1.
Therefore (H\_new,P,theta\_new) is an actual arithmetic finite
Heisenberg pair, and its actual irreducible carrier E has dimension2.

All representations here have finite image on the compact group P: theta
has open kernel, the induction has finite dimension, and an open normal
core of that kernel acts trivially. Thus no infinite projective-image
issue is being suppressed.

This proof does not require rescaling the actual matrix g inside GL2(F).
Its represented square theta(g²) can be an arbitrary retained scalar.
Rescaling X is only allowed later in a C× pushout if that is the
attachment contract.

\hypertarget{rank2-explicit-embedding-into-the-same-fixed-candidate}{%
\subsection{5. Explicit embedding into the same fixed
candidate}\label{rank2-explicit-embedding-into-the-same-fixed-candidate}}

Realize Pi as functions F:G→C with

\begin{Verbatim}[fontsize=\small,breaklines=true,breakanywhere=true]
F(nx)=Lambda(n)F(x),
\end{Verbatim}

compact support modulo N, and right translation by G. For f∈E define
i(f) by

\begin{Verbatim}[fontsize=\small,breaklines=true,breakanywhere=true]
i(f)(np)=Lambda(n)f(p),       n∈N, p∈P,
\end{Verbatim}

and zero off NP. This is well-defined because N∩P=J and
Lambda\textbar J=theta. The set NP is a finite union of N-cosets, by
{[}P:J{]}=2, so the support is compact modulo N. The map i is
P-equivariant and injective.

Its image is exactly the two-dimensional P-stable subspace of Pi
consisting of inducing functions supported on NP. Thus the enlarged
carrier is already inside the given Pi. No new candidate is chosen, and
all noncompact values of Lambda remain fixed.

Since Pi is irreducible and i(E) is nonzero, its G-span is Pi. More
precisely, there is the canonical G-equivariant quotient

\begin{Verbatim}[fontsize=\small,breaklines=true,breakanywhere=true]
Q_E:c-Ind_P^G E→Pi,
Q_E(F)=sum_(Px∈P\G) Pi(x^−1)i(F(x)).                 (7)
\end{Verbatim}

The sum is finite and independent of representatives. If x is replaced
by px, then F(px)=E(p)F(x) and P-equivariance of i cancels the
additional Pi(p\^{}−1). Formula (7) is nonzero and hence onto. The next
section identifies its kernel exactly.

\hypertarget{rank2-exact-recovery-by-the-noncompact-hecke-eigenvalue}{%
\subsection{6. Exact recovery by the noncompact Hecke
eigenvalue}\label{rank2-exact-recovery-by-the-noncompact-hecke-eigenvalue}}

Let X=c-Ind\_J\^{}G theta, in the same left-covariant/right-translation
convention. Induction in stages gives an explicit isomorphism

\begin{Verbatim}[fontsize=\small,breaklines=true,breakanywhere=true]
c-Ind_P^G E → X,
F ↦ [x↦F(x)(1)].                                     (8)
\end{Verbatim}

The inverse sends f∈X to the E-valued function F(x)(p)=f(px). Both
covariance statements follow directly from the definitions, and the
finite index {[}P:J{]} ensures the support conditions agree.

Fix a uniformizer varpi=varpi\_K. Because K× normalizes theta, the
operator

\begin{Verbatim}[fontsize=\small,breaklines=true,breakanywhere=true]
(Tf)(x)=f(varpi x)                                   (9)
\end{Verbatim}

is a well-defined invertible G-equivariant operator on X. Put
s=Lambda(varpi), which is nonzero. Since k\_K=F2 and J contains
U\_K\^{}1,

\begin{Verbatim}[fontsize=\small,breaklines=true,breakanywhere=true]
N/J≅K×/U_K^1≅Z,
\end{Verbatim}

with generator varpi. Define

\begin{Verbatim}[fontsize=\small,breaklines=true,breakanywhere=true]
Q:X→Pi,
(Qf)(x)=sum_(ell∈Z) s^(−ell) f(varpi^ell x).          (10)
\end{Verbatim}

For each x the sum is finite, since f has compact support modulo the
open subgroup J. Equation (10) satisfies

\begin{Verbatim}[fontsize=\small,breaklines=true,breakanywhere=true]
(Qf)(varpi x)=s(Qf)(x),
(Qf)(jx)=theta(j)(Qf)(x),       j∈J,
\end{Verbatim}

and therefore the required N-covariance. Its support is compact modulo
N. It commutes with G-right translation and satisfies

\begin{Verbatim}[fontsize=\small,breaklines=true,breakanywhere=true]
QT=sQ.                                               (11)
\end{Verbatim}

It is onto: a function in Pi has finite support modulo N, and on each
occupied N-coset its value can be lifted to a function supported on one
J-coset. Extend those finitely many choices using theta.

To compute the kernel, decompose X as the algebraic direct sum of its
subspaces supported on the left N-cosets in G. Each such summand is the
space of finitely supported sequences on N/J≅Z, with the fixed theta
covariance. The operator T is a bilateral shift. On that summand, (10)
is evaluation of its Laurent polynomial at the nonzero scalar s, up to
the harmless convention of replacing the Laurent variable by its
inverse. The kernel is therefore the principal ideal generated by T−s.
The elementary identity

\begin{Verbatim}[fontsize=\small,breaklines=true,breakanywhere=true]
ker(C[z,z^−1]→C, z↦s)=(z−s)
\end{Verbatim}

proves, summand by summand, that

\begin{Verbatim}[fontsize=\small,breaklines=true,breakanywhere=true]
ker Q=(T−s)X.
\end{Verbatim}

All vectors occupy finitely many such summands, so no completion or
closure is needed. Thus

\begin{Verbatim}[fontsize=\small,breaklines=true,breakanywhere=true]
Pi ≅ X/(T−s)X ≅ (c-Ind_P^G E)/(T−s)(c-Ind_P^G E),    (12)
\end{Verbatim}

where T is transported through (8).

Under (8), formula (10) agrees with (7): on the standard copy of E
supported on P, it gives exactly i, and G-equivariance then gives
equality everywhere. Thus (12) identifies the kernel of the canonical
recovery map, not merely the kernel of an unrelated surjection.

The same description can be phrased as evaluation of the type Hecke
algebra at the character determined by Lambda. Its intertwining support
is N, each supported intertwining space is one-dimensional, and the
retained extension Lambda normalizes the Laurent generator. The explicit
formulas above avoid any ambiguity about convolution orientation.

The scalar s is essential. E alone, being compact data, does not
determine it. Changing s can change the unramified normalization of the
candidate. The construction retains the original Lambda and s, so it
recovers the exact original Pi rather than an unspecified member of its
inertial class.

\hypertarget{rank2-full-scalar-pushout-attachment-and-remaining-marking-obligations}{%
\subsection{7. Full scalar-pushout attachment and remaining marking
obligations}\label{rank2-full-scalar-pushout-attachment-and-remaining-marking-obligations}}

Form the C× pushout of the actual pair (H\_new,P,theta\_new). It is a
Heisenberg extension of the actual symplectic plane P/H\_new. Any
supplied source plane with the same chosen alternating form has a
center-preserving pushout isomorphism into it. Every source
full-character value is then preserved in the center, as in the retained
full-character attachment theorem. Stone--von Neumann identifies its
source carrier with the actual two-dimensional E.

Together with the embedding i and exact quotient (12), this provides an
actual arithmetic carrier and recovery inside the same fixed candidate.
It closes the carrier-capacity failure of the old rank-two eta by
changing the compact arithmetic pair, not by increasing the old eta or
discarding its dimension calculation.

At p=2, a C× pushout does not automatically identify a separately marked
native mu\_2 extension or quadratic refinement. The source finite
central extension can be retained through its chosen embedding, while a
required comparison with the target's preselected square function
remains an additional question. No arbitrary native Arf invariant is
asserted here.

Nontrivial stabilizer operators, transport graphs, and original H2's
requirement that a target be a particular BK type-side quotient are
likewise separate. In particular, the theorem proves that the new pair
is an actual arithmetic finite Heisenberg pair, but does not rename it
as the unchanged maximal-simple-type pair. No new hypothesis is silently
added to original Theorem16.2.

\hypertarget{rank2-scope-smallest-boundary-and-sources}{%
\subsection{8. Scope, smallest boundary, and
sources}\label{rank2-scope-smallest-boundary-and-sources}}

The inequality q\textgreater=3 matters. If q=1, then m=0 and the
displayed enlargement is not the group used above. Over residue field F2
the complete Iwahori unit group equals U\_A\^{}1, and a naive level-zero
replacement does not give the asserted index-two enlargement. No q=1
claim is made. The source levels q=3h with positive odd h satisfy the
required q\textgreater=3.

The proof uses the actual one-dimensional Lambda because the maximal
parameter field has full degree and residue field F2. It does not apply
the same scalar argument automatically to a higher-rank inducing
representation.

Primary BK passages: 3.1.7--8 (odd minimal H/J equality); 3.2.1--2 (deep
restriction and K× normalization); 3.3.2 and 5.1.1/5.1.8 (intertwining);
and 6.2.1--5 (the fixed maximal-type extension and compact induction).
The subgroup, matrix, Mackey-support, and Laurent-evaluation
calculations are proved explicitly above. No external correspondence
theorem is used in this enlargement.

Retained primary text:

\begin{Verbatim}[fontsize=\small,breaklines=true,breakanywhere=true]
closure_recovery/extracted/Paper_XI_Continuation158_General_Sectors_and_Arit
hmetic_Constraints/scratch158/information158/input_cache/BK_source.txt
SHA256 d69db0827256707a34cd2562af68b9de11fc889ab3f43535ec47b620d58ca9c8
\end{Verbatim}

The previously reviewed negative statement concerned eta of the original
BK \(J^1/H^1\). That quotient remains zero. The new positive
statement concerns a different, explicitly constructed arithmetic pair
inside the same ambient representation, with exact
noncompact-normalization-aware recovery.
